\begin{abstract}
Service robots in crowded pedestrian environments re-decide, at every control cycle, which side --- left or right --- to pass a person on; when this re-decision is greedy, the passing direction oscillates across cycles in ambiguous configurations, or the robot freezes. Memoryless local planners keep no record of the previous passing decision. Topology-aware planners bias selection toward the previously chosen class through soft costs, selection hysteresis, or consistency weights, and opinion-dynamics methods make the passing side an explicit, stateful decision with deadlock-breaking guarantees; to our knowledge, none of these states a bound on how often a discrete passing decision may flip. We introduce \emph{\sname} (COMmitment-based PASSing for Social navigation), a decision-layer planner that treats the \emph{temporal consistency} of passing decisions as a first-class design objective. The passing relationship is represented as a topological \Class\ bound to a person's track ID, so that the previous decision can be re-identified across cycles (a design contract that the current implementation only partly realizes), and safety is layered lexicographically, as in standard priority arbitration, on top of a commitment switching rule that folds the switching margin, dwell, and point of no return into a single \emph{leaky evidence accumulation} equation --- a rectified, leaky variant of Page's one-sided CUSUM statistic with a progress-dependent threshold. We state five properties (P1--P5) --- anti-oscillation, switching-rate bound, safety dominance, maneuver completion, and conditional finite-window escalation to HOLD --- obtained by applying standard dwell-time and CUSUM/queueing arguments to this rule; they are not new mathematics. The anti-oscillation guarantee we claim is \emph{deterministic and pathwise}: for fixed sampling period $dt$, at least $N_{\min} = \lceil E_0/((D_{\max}-\Dfloor)\,dt)\rceil$ control cycles separate consecutive discretionary switches (P2, an elementary bounded-increment bound of dwell-time type), assuming only the range bound $D_k \le D_{\max}$ that the $[0,1]$ cost normalization supplies by construction, together with the declared non-vacuity condition $E_0 > (D_{\max}-\Dfloor)\,dt$ ($N_{\min} \ge 2$), which the default knobs satisfy ($N_{\min} = 3$ cycles in the worst case $D_{\max} = \sum_X w_X$, and $7$ cycles when the advantage is bounded by a single normalized term). P2 bounds the switching rate from above only; a companion response-time bound gives, under a sustained challenger advantage, the finite number of cycles within which a discretionary switch occurs, under persistent selection, uninterrupted updates, and sufficient clip and equilibrium headroom. P1 refines this, for i.i.d.\ challenger advantages, to an exponentially small per-cycle switching probability $e^{-\theta^* E_0}$ by dominating the rule with the Lindley (CUSUM) recursion and applying Kingman's bound, where $\theta^*$ is the exact Lundberg root of the increment law. In an offline harness that drives the actual decision core (5 scenarios $\times$ 50 seeds), an immediate-argmin comparator produces about 98 decision changes per encounter under ambiguous ties, whereas \sname\ produces at most one decision change in each tested encounter; cyclic latency at $K=3$ (8 classes) is p99 $58.6\,\mu$s, max $651.3\,\mu$s --- about $1.30\%$ of the 20~Hz budget (50~ms) --- in a historical decision-core-only measurement; a decision-stream endpoint-suffix statistic is reported that is not a validated legibility measure; and an offline progress-rate transition-blocking interval $v_{\text{lat}} \in (0.20, 0.35)$~m/s is reported, not a physical freezing threshold. Physical performance metrics, closed-loop evaluation, external baselines (including a previous-class consistency-weight baseline), and user studies are not measured; they are specified as a planned-evaluation protocol. Code: \url{https://github.com/kjungmo/compass}.
\end{abstract}
